\documentclass[11pt,a4paper]{article}
\usepackage[utf8]{inputenc}
\usepackage[T1]{fontenc}
\usepackage[slovak]{babel}
\usepackage[margin=2.5cm]{geometry}
\usepackage{amsmath,amssymb}
\usepackage{graphicx}
\usepackage{booktabs}
\usepackage{hyperref}

\title{Replikácia iterovaného greedy algoritmu pre flow shop\\s chýbajúcimi operáciami}
\author{Lukáš \and Dominik \and Jano}
\date{\today}

\begin{document}
\maketitle

% Rozsah: 6–10 strán A4 bez príloh. Každá tabuľka a graf musí mať popis a byť interpretovaná v texte.

\begin{abstract}
% TODO
\end{abstract}

\section{Opis a význam riešeného problému}
% TODO: flow shop, chýbajúce operácie, mass customization; \cite{RittRossit2024}.

\section{Formulácia optimalizačnej úlohy}
% TODO: F | skip | \sum C_j; typ problému (min., diskrétny, jednokriteriálny,
% deterministický, statický, s obmedzeniami, white-box); rozhodovacie premenné;
% priestor riešení a reprezentácia (permutácia / pseudo-joby \cite{BenavidesRitt2018});
% obmedzenia; účelová funkcia; vyhodnocovanie kvality.

\section{Metóda autorov článku}
% TODO: dvojfázový IGA \cite{RittRossit2024}, \cite{RuizStutzle2007},
% vkladanie s anticipáciou/oneskorením \cite{BenavidesRitt2015}, porovnávaná 2SA \cite{HennebergNeufeld2016}.

\section{Naša implementácia}
% TODO

\section{Rozdiely oproti pôvodnému článku}
% TODO: jazyk (Python + numba vs. C++), hardvér, inštancie, počet seedov, prípadné zmenšenie rozsahu.

\section{Experimentálna metodológia}
% TODO: inštancie, parametre, rozpočet 15000/n iterácií, seedy 1–10, metriky.

\section{Dosiahnuté výsledky}
% TODO

\section{Porovnanie s výsledkami autorov}
% TODO: Tab. 3, 4, 6 z článku (data/paper_results).

\section{Diskusia a možné zdroje rozdielov}
% TODO

\section{Obmedzenia experimentov}
% TODO

\section{Záver}
% TODO

\bibliographystyle{plain}
\bibliography{../CITATION}

\end{document}
